% critique.tex — f176v, vv. 1169–1210 : commentaire critique.
% Parties textuelles du commentaire de l'auteur (travaux antérieurs, Harmattan 2023), texte et notes
% conservés ; répartition, harmonisation et passages signalés : repartition.md.
\cc{1174}{πρῶτον μὲν οὖν δίελθε κἀνακόλπασον.}{(Que le lecteur curieux aille faire son miel de la note qu'a consacrée Jacques Moreau à ce vers\nc{«~Note sur le vers 1174 des \italique{Thesmophories} d'Aristophane~», \italique{L'antiquité classique}, tome 7, fasc. 2 (1938) pp. 227-231. À peine âgé de vingt ans, le jeune étudiant d'alors fait montre d'une érudition et d'un esprit critique peu communs ; il faut dire qu'il avait pour professeur à l'Université de Bruxelles Henri Grégoire -- le successeur d'Alphonse Willems -- et dont l'assistant n'était autre que Roger Goossens ! bouillon intellectuel dont on imagine sans peine l'extraordinaire émulation qui devait en résulter parmi les élèves.}.)
Le \italique{Ravennas} porte au vers 1174 la leçon : πρῶτον μὲν οὖν δίελθε κἀνακ\textbf{ό}λπασον.
Le verbe ἀνακολπάζειν est un hapax, certes, mais dont le sens se laisserait entendre -- s'il n'y avait lieu comme nous l'allons voir de le soupçonner : ὁ κόλπος désigne au propre un \italique{pli}, \italique{creux} ; aussi bien le pli d'un vêtement -- notamment au niveau de la ceinture -- que, chez la mère ou la nourrice, le \italique{genou} ou le \italique{sein}. On trouve chez Homère l'expression κόλπον ἀνιεμένη\nc{\italique{Iliade}, XX, 80. Et l'explication d'Hésychius (κ 3393) : κόλπον ἀνιεμένη· τὸ κατὰ τὸ στῆθος μέρος ἀνέλκουσα, ἵνα τὸν μαστὸν δείξῃ.} lorsqu'Hécube fait tomber le haut de son vêtement et découvre son sein.
Il n'y aurait rien d'étonnant à ce qu'ἀνακολπάζειν ait pu signifier \italique{retrousser un vêtement}\nc{Ou, à la rigueur, «~tenir haut les seins~» comme le traduit Germaine Prudhommeau.}.
Pour définir les ζειραί, ces vêtements qu'on \italique{resserre au niveau de la ceinture}, le rédacteur anonyme de l'\latin{Etymologicum Magnum} use du verbe ἀνακολπόω :\citComm{ζειραί: χιτῶνες ἀνακεκολπωμένοι\nc{†ἀνακεκολλαμένοι† dans divers lexiques, notamment le lexique platonicien de Timée le Sophiste ou celui de Photius. La correction est-elle du fait du rédacteur de l'\latin{Etymologicum Magnum} ?}, ἢ ἀνάκωλοι.}Si le \italique{Ravennas} ne comportait en sus du texte au vers 1174 une scholie marginale venant renverser tout notre fragile édifice, ἀνακολπάζειν ferait partie de ces nombreux hapax dont le sens n'a été inféré qu'à grand renfort d'imagination sur de maigres indices étymologiques.\\
Cependant un ancien commentateur eut l'heureuse idée de gloser la leçon que son manuscrit donnait :\citComm{σ\douteux{η}[ * * * ] τοῦτο. ση[μαίνει] δὲ τὸ `ἁβρῶς βαδίζειν'.}Exit donc la leçon ἀνακολπάζειν, sans rapport avec le \italique{déplacement langoureux}\nc{Ou \italique{lascif}, dira-t-on, au sens que le \nieme{XXI} siècle prête à ce mot ; c'est-à-dire avec \italique{langueur}, un \italique{manque de vigueur} qui trahit un fort penchant pour la sensualité. Est-ce l'influence de \italique{las} qui aura changé radicalement le sens originel de \italique{lascif} (\italique{lubrique}, \italique{folâtre} ; sens que rapportent encore exclusivement nos dictionnaires) ?} dont parle le scholiaste, nullement étayée et d'ailleurs inepte : la danseuse court-vêtue n'a aucun vêtement pesant à retrousser ; au moment d'enflammer le Scythe, Euripide lui fera sauter les préliminaires en lui ordonnant à brûle-pourpoint de retirer tout bonnement son léger ἱμάτιον.\\
Hermann\nc{Dans sa recension des \italique{Thesmophories} éditées par Enger en 1844 (\italique{Aristophanis Comoediae cum scholiis}), in \italique{Zeitschrift für altertumswissenschaft} (1845), coll. 905-920. Article réédité par Fritzsche dans les \italique{Godofredi Hermanni Opuscula}, vol. 8 ; Hermann y discute de notre vers à la page 307.} proposa de lire le verbe ἀνακαλπάζειν\nc{Au vers 1201, un `α' est lu fautivement de la même manière en `ο'.}. 
Son hypothèse a été reprise depuis lors par toutes les éditions ; elle s'appuyait effectivement sur un indice de taille : à l'adresse κ 623 de la \italique{Souda}, l'on retrouve notre glose précédée du lemme καλπάζειν :\citComm{καλπάζειν· τὸ ἁβρῶς βαδίζειν.}Le verbe est formé à partir de ἡ κάλπη, \italique{le trot} (ou le \italique{galop} ?), qui désigne une sorte de course aux jeux olympiques, où le cavalier, en approchant du but, devait prendre le trot, sauter à terre et suivre son cheval.
Quant au composé, ἀνακαλπάζειν, Photius (α 1514) nous en a conservé une trace en recopiant la glose d'un grammairien versé dans le \italique{Bon usage} attique :\citComm{ἀνακαλπάζει· τινὲς μὲν ὡς οὐ δόκιμον ἐφυλάξαντο τὴν φωνήν, Αἰσχύλος δὲ ἐχρήσατο Μυσοῖς \{ὡς δόκιμον\}· λέγει γάρ `εἶδον καλπάζοντας ἐν αἰχμαῖς'. ὁμοίως Σοφοκλῆς καὶ Ἀριστοφάνης καὶ Πλάτων καὶ ἕτεροι.\\
ἀνακαλπάζει· certains défendirent [l'usage de] ce mot, comme non attesté. Eschyle pourtant l'utilisa dans ses \italique{Mysiens}. Il dit en effet : «~Je les vis trotter au milieu des lances~». De même Sophocle, Aristophane, Platon et d'autres.}C'est bien le sens de \italique{trotter}, \italique{galoper}, c'est-à-dire \italique{courir} qu'il faut donner au verbe <ἀνα>καλπάζειν, témoin le fragment d'Eschyle cité par Photius\nc{$=$ fr. 144a Radt.} : il y fait allusion à une locution proverbiale\nc{Transmis par les parémiographes Zénobius (III, 13), Suidas (δ 1186 ; ο 414) et Érasme (2014). Il n'est pas impossible que le proverbe soit lui-même tiré de la danse, et d'une éventuelle \italique{danse des épées}. Sur l'utilisation d'épées pour certaines acrobaties périlleuses, voir \italique{La danse grecque antique} de Maurice Emmanuel (1896), fig. 556 et 557 ou l'article «~Histoire anthropologique des danses acrobatiques~» de	Myriam Peignist, \italique{Corps} (2009/2), n° 7, pp. 29-38. La danseuse se livre-t-elle à ce genre d'acrobaties guerrières pour attirer l'attention du Scythe ? cela paraît peu vraisemblable ; il est préférable d'interpréter le fragment d'Eschyle de façon imagée et ne pas rattacher directement le verbe καλπάζειν à une quelconque danse des épées.
Un proverbe analogue et que l'on connaît en français, \italique{être sur le fil du rasoir}, apparaît déjà chez Homère (\italique{Iliade} X, 173).} :\citComm{`δι᾽ ὀξείας δραμεῖν'· ἐπὶ τῶν διακινδυνευόντων. ὀξεῖαν γὰρ λέγουσι τὴν λόγχην\\
«~Courir en traversant les pointes~» : à propos de ceux qui prennent des risques ; on appelle pointe la lance.}Dans le passage d'Eschyle, καλπάζοντας est un permutant métaphorique -- donc synonyme -- de δραμόντας.\\
Willems \& nombre d'exégètes voulurent rattacher <ἀνα>καλπάζειν à un terme technique de danse : \italique{se déplacer} ou \italique{danser sur les pointes}. Le sens leur a-t-il été suggéré par l'image du cheval au trot ? Dans le célèbre roman de Longus (II, 37), lorsque Daphnis imite le dieu Pan avec ses sabots, il se met à courir \citComm{ἐπ' ἄκρων τῶν ὀνύχων.}Là est le point : car l'on sait le type d'entraînement auquel se soumet la danseuse : au vers 1180, le Scythe qui la voit si \italique{légère} se mettre à \italique{sauter} («~comme une puce~», dixit le Scythe...), il s'exclame : ὡς ἐλαπρός, ὥσπερ ψύλλο κατὰ τὸ κῴδιο.
Moreau avait donc raison de souligner que la glose présentant notre danseuse aguichant l'Archer \italique{mollement} (ἁβρῶς βαδίζειν) ne tenait pas ; mais l'erreur venait de la glose et non de la correction d'Hermann : on eut le tort de vouloir conserver à tout prix cette glose telle que rapportée par le \italique{Ravennas} et Suidas ; elle n'était pourtant en conformité ni avec l'action ni avec le type d'entraînement -- hyperactif -- de la danseuse.
Le lexique de Jean Zonaras\nc{Édité par Tittmann sous le titre de \italique{Iohannis Zonarae Lexicon} (1808), p. 1171. Ce lexique de la fin du \nieme{XII} ou début du \nieme{XIII} siècle, s'il faut en croire Klaus Alpers, serait plutôt à attribuer à un certain Nicéphore. Le Pseudo-Zonaras a puisé abondamment dans la \italique{Souda} (ou des sources très voisines), notamment parmi les adresses de la lettre Α à Κ et les gloses touchant à Aristophane. Notre cas est intéressant, puisqu'il est peu vraisemblable que la transposition de β et κ soit une erreur ou correction du copiste ; il faut donc en conclure que si le Pseudo-Zonaras est allé piocher sa glose dans la \italique{Souda}, il disposait d'un meilleur manuscrit que ceux qui nous sont parvenus. Klaus Alpers a consacré un article très riche au lexique du Pseudo-Zonaras dans la monumentale \italique{Realencyclopädie der classischen Altertumswissenschaft}, 10/1 (1972), à l'entrée \allemand{'Zonarae' Lexikon}, coll. 732-763. Il le résume dans \italique{Der Neue Pauly}, 12/2 (2003), col. 831. Pour une étude plus générale de l'ensemble du corpus des anciens lexiques en langue grecque, l'on consultera la synthèse du même auteur dans l'\italique{Historisches Wörterbuch der Rhetorik}, 5 (2001) à l'entrée \allemand{Lexikographie} B, coll. 194-210.} nous en livre la bonne version (simple substitution de ἄ\textbf{κ}ρως à ἁ\textbf{β}ρῶς) :\citComm{καλπάζειν· τὸ ἄκρως βαδίζειν.}Willems avait donc vu juste : l'exercice de la danseuse devait évoquer la course de la κάλπη : commencer par un \italique{trot}, c'est-à-dire par un déplacement en \italique{pointes} (τὸ ἄκρως βαδίζειν) qui s'accélérait pour terminer par un saut ; ce qu'Euripide désigne par le verbe ἀνακαλπάζειν.\\
Cette interprétation nous est confirmée par Eustathe (\italique{Commentaire à l'Iliade}, XX, v. 228 p. 1206, à propos du mot `σκιρτῷεν') :\citComm{ἔστι δὲ σκιρτᾶν οὐ τὸ ἁπλῶς θέειν, ἀλλὰ τὸ ἀνέτως καὶ δίχα ἐπιταγῆς καὶ πηδητικῶς καὶ οἷον ὀρχηστικῶς \textbf{κατὰ τοὺς καλπάζοντας}.\\
«~Bondir~» (σκιρτᾶν) n'est pas le simple fait de courir, mais de le faire librement, sans ordre, fougueusement, ainsi que le font les \textbf{καλπάζοντες} en dansant.}Terminons par Hésychius (κ 519)\nc{Contrairement à l'usage de mon commentaire, je n'ai réservé Hésychius que pour la fin de cette note. La raison en est que la glose de l'unicodex -- un \italique{Marcianus} du \nieme{XV} siècle -- est légèrement corrompue et donc sujette à caution :\citComm{καλπάγει ὀξυπόδης}(Le lexique de Cyrille ajoute : σακκιζει (sic, immo βαδίζει).)}, qui ne le comprenait pas autrement :\citComm{καλπάζει· ὀξυποδεῖ.}ὀξυποδέω signifie \italique{accélérer le pas}\nc{Cf. le passage de l'\italique{Oreste} d'Euripide, vers 1550, où Ménélas furieux accourt ὀξύπους ($=$ ταχύπους). Ce sens est celui rapporté avec raison par le \LSJ. Le mot latin \latin{addensator}, \italique{celui qui rend plus serré} (en parlant du foulon et de son pas), est glosé dans certains manuscrits gréco-latins ὀξυποδητής.}. Notons aussi l'autre glose d'Hésychius (σ 822) :\citComm{σκαλ\{α\}πάζειν· ῥεμβωδῶς βαδίζειν.\\
σκαλπάζειν : se déplacer en tournoyant.}}
\cc{1176}{κωμοτίς ἀνεγεῖρί μοι;}{Variante de lecture du scholiaste : διχῶς κωμοδρίαν.}
\cc{1181}{φέρε θοἰμάτιον ἄνωθεν.}{Sur le texte et le sens de ce vers 1181, voir la note 28 d'Alphonse Willems, reproduite ci-après.}
\cc{1185–93}{Papyrus d'Oxyrhynque n° 3840.}{Un mince débris contenant quelques lettres de notre passage fut déniché parmi les papyrus d'Oxyrhynque\nc{\italique{The Oxyrhynchus Papyri}, vol. LVI (1989), n°3840. Étudié par Peter J. Parsons, aidé de Colin Austin -- déjà en train de préparer son édition des \italique{Thesmophories}. La transcription qu'en fit Parsons et son analyse sont recopiées en annexe. Voir aussi le papyrus d'Oxyrhynque n° 5132 (\italique{The Oxyrhynchus Papyri} LXXVIII) pour un semblant d'\italique{addendum} : la présence de doubles points `$ : $', comme au vers 1190, pour indiquer un changement de locuteur est confirmée par un autre fragment du même papyrus au vers 1209.}. Il y a néanmoins quelques informations à en tirer : il confirme la conjecture d'Enger au vers 1185 (γογγυλί à la place de γογγύλῃ), ainsi que μέλι au vers 1192. Parsons note la présence d'une ligne supplémentaire après le vers 1187, là où le \italique{Ravennas} possède l'une de ses rares didascalies\nc{Pour notre pièce où les scholies très anciennes et érudites se font rares, on ne relève d'autres didascalies qu'aux vers 100, 129 et 276.} : ἀνακύπτη· καὶ παρακύπτι· ἀπεψωλημένος· (même jeu de scène avec l'olisbos qu'aux vers 643-648).}
\cc{1197–216}{ἀλλὰ τὸ συμβήνην λαβέ $[$ ... $]$ ὀρτῶς δὲ συβήνη 'στί· καταβηνῆσι γάρ.}{Pour le sens de συμβήνη et pour l'action, voir le commentaire de la pièce.\\
Et je me garderai bien de biffer le μ de la leçon du manuscrit, συ\textbf{μ}βήνην, qui aura survécu à plus de mille ans de tradition manuscrite mais pas à la griffe implacable des éditeurs d'Aristophane.\\
Le début du vers suivant, altéré, ἔπειτα †κ\inRas{ω}{ο}μίζεις αὐτοῖς†, se laisse lire sans difficulté :\citComm{ἔπειτα κομίζι ($=$ κομίζω) σ(οὶ) αὐτίκ(α)\\
Ensuite je te l'apporte $[$la drachme$]$ sur-le-champ.}La dernière version du manuscrit attribue, il est vrai, cette réplique à Euripide. Ces indications marginales de personnages ne doivent pas nous arrêter : elles sont récentes et même pour certaines du \nieme{XVI} siècle\nc{Cf. \italique{supra}, v. 1059 n. 5.} (c'est le cas de cette annotation, «~γϋ.~», que porte le manuscrit (la même qu'au vers 1189) et mis pour γυνή -- c'est-à-dire la \italique{vieille}, γραῦς ; c'est-à-dire Euripide. Von Velsen ne l'a pas intégrée à son apparat critique et avait donc reconnu l'écriture d'Euphrosynus Boninus).\\
Mon émendation appelle quelque justification : l'hypothétique confusion αυτικ / αυτοις peut aisément s'expliquer en constatant la ressemblance entre le digramme `ις' et le `κ'\nc{Il faudrait alors supposer l'existence d'une forme intermédiaire αυτιις ou αυτις, bien vite corrigée en αυτοις. Une forme intermédiaire, ou un \italique{chaînon manquant}, pour reprendre le jargon des darwiniens. Ne poussons donc pas plus loin la «~preuve~» afin d'éviter l'écueil de ces demi-savants, capables de tout prouver et même son contraire !}. Continuons.\\
Pour l'action et le trait d'esprit du vers 1215 (ὀρτῶς δὲ συμβήνη ᾽στί· καταβηνῆσι γάρ), voir le commentaire de la pièce. Nul besoin de retailler le vers : καταβαίνω est le terme consacré lorsqu'un acteur descend de la σκηνή vers l'ὀρχήστρα, d'autant plus adapté quand la scène se situe sur un lieu en hauteur comme la colline de la Pnyx\nc{Le Scythe s'interrompt quelques instants de parler, pause soulignée par l'interjection \italique{extra metrum} οἴμοι qu'il lance une fois descendu : il se demande vers laquelle des deux entrées de l'orchestre se tourner (τί δρᾶσι; ποῖ τὸ γρᾴδι᾽;)} ; cf. les \italique{Guêpes} (vers 1514) ou l'\italique{Assemblée des femmes} (vers 1152)\nc{\italique{Confer} surtout la note d'Alphonse Willems au vers 1341 des \italique{Guêpes}. Lisez Willems ! aucune excuse : la note est reproduite en annexe.}.}
\cc{1203–11}{Papyrus d'Oxyrhynque n° 4935 fr. 2 \& n° 5132.}{Deux fragments d'un même papyrus d'Oxyrhynque\nc{N° 4935 fr. 2, in \italique{The Oxyrhynchus Papyri}, LXXIII (2009), ed. A. Benaissa. N° 5132, \italique{The Oxyrhynchus Papyri}, LXXVIII (2012), ed. W. B. Henry.} contiennent quelques lettres de ce passage. (Voir en annexe.) Ils n'apportent aucune leçon originale. Au vers 1208, l'espace manque pour être compatible avec la leçon \italique{ante correctionem} λέλυσον du \italique{Ravennas} ; l'état du fragment ne permet pas pour autant de trancher entre les corrections λέλυσο (\Rpc) et λέλυσαι (Bentley).}
